% Physical source leaf 019; retained native image authority.
\SourceLeaf{19}
\begin{center}
{\large\bfseries\textls[120]{SECTIO SECUNDA}}\par\bigskip
{\small\bfseries DE}\par\bigskip
{\large\bfseries CONGRUENTIIS PRIMI GRADUS.}
\end{center}
\begin{center}\rule{0.09\linewidth}{0.3pt}\end{center}
\begin{center}\emph{Theoremata praeliminaria de numeris primis, factoribus etc.}\end{center}
\Article{13}
\textsc{Theorema}. \emph{Productum e duobus numeris positivis numero primo dato minoribus per hunc primum dividi nequit.}

Sit $p$ primus, et $a$ positivus $<p$: tum nullus numerus positivus $b$ ipso $p$ minor dabitur, ita ut sit $ab\equiv 0\ (\mathrm{mod.}\ p)$.

\emph{Dem.} Si quis neget, supponamus dari numeros $b, c, d$ etc. omnes $<p$, ita ut $ab\equiv 0$, $ac\equiv 0$, $ad\equiv 0$ etc. (mod. $p$). Sit omnium minimus $b$, ita ut omnes numeri ipso $b$ minores hac proprietate sint destituti. Manifesto erit $b>1$: si enim $b=1$, foret $ab=a<p$ (\emph{hyp.}), adeoque per $p$ non divisibilis. Quare $p$ tamquam primus per $b$ dividi non poterit, sed inter duo ipsius $b$ multipla proxima $mb$ et $(m+1)b$ cadet. Sit $p-mb=b'$, eritque $b'$ numerus positivus et $<b$. Iam quia supposuimus.\ $ab\equiv 0\ (\mathrm{mod.}\ p)$, habebitur quoque $mab\equiv 0$ (art. 7), et hinc.\ subtrahendo ab $ap\equiv 0$, erit $a(p-mb)=ab'\equiv 0$; i. e. $b'$ inter numeros $b, c, d$ etc. referendus, licet minimo eorum $b$ sit minor. \emph{Q.~E.~A.}

\Article{14}
\emph{Si nec $a$ nec $b$ per numerum primum $p$ dividi potest: etiam productum $ab$ per $p$ dividi non poterit.}

% Physical source leaf 020; retained native image authority.
\SourceLeaf{20}
Sint numerorum $a, b$, secundum modulum $p$ residua minima positiva $\alpha, \beta$, quorum neutrum erit $0$ (\emph{hyp.}) Iam si esset $ab\equiv 0$ (mod. $p$), foret quoque, propter $ab\equiv\alpha\beta$, $\alpha\beta\equiv 0$, quod cum theoremate praec. consistere nequit.

Huius theorematis demonstratio iam ab Euclide tradita, \emph{El.} VII. 32. Nos tamen omittere eam noluimus, tum quod recentiorum complures seu ratiocinia vaga pro demonstratione venditaverunt, seu theorema omnino praeterierunt, tum quod indoles methodi hic adhibitae, qua infra ad multo reconditiora enodanda utemur, e casu simpliciori facilius deprehendi poterit.

\Article{15}
\emph{Si nullus numerorum $a, b, c, d$ etc. per numerum primum $p$ dividi potest, etiam productum $abcd$ etc. per $p$ dividi non poterit.}

Secundum artic. praec. $ab$ per $p$ dividi nequit; ergo etiam $abc$; hinc $abcd$ etc.

\Article{16}
\textsc{Theorema}. \emph{Numerus compositus quicunque unico tantum modo in factores primos resolvi potest.}

\emph{Dem.} Quemvis numerum compositum in factores primos resolvi posse, ex elementis constat, sed pluribus modis diversis fieri hoc non posse perperam plerumque supponitur tacite. Fingamus numerum compositum $A$, qui sit $=a^\alpha b^\beta c^\gamma$ etc., designantibus $a, b, c$ etc. numeros primos inaequales, alio adhuc modo in factores primos esse resolubilem. Primo manifestum est.\ in secundum hoc factorum systema alios primos quam $a, b, c$ etc. ingredi non posse, quum quicunque alius primus numerum $A$ ex his compositum metiri nequeat. Similiter etiam in secundo hoc factorum systemate nullus primorum $a, b, c$ etc. deesse potest, quippe qui alias ipsum $A$ non metiretur (art. praec.). Quare hae binae in factores resolutiones in eo tantummodo differre possunt, quod in altera aliquis primus pluries quam in altera habeatur. Sit talis primus $p$, qui in altera resolutione $m$, in altera vero $n$ vicibus occurrat, sitque $m>n$: Iam deleatur ex utroque systemate factor $p$, $n$ vicibus, quo fiet ut in altero adhuc $m-n$ vicibus remaneat, ex altero vero omnino abierit. I. e. numeri $\frac{A}{p^n}$ duae in factores resolutiones habentur, quarum altera a factore $p$ prorsus libera, altera vero $m-n$ vicibus eum continet, contra ea quae modo demonstravimus.

% Physical source leaf 021; retained native image authority.
\SourceLeaf{21}
\Article{17}
Si itaque numerus compositus $A$ est productum ex $B, C, D$ etc., patet, inter factores primos numerorum $B, C, D$ etc. alios esse non posse, quam qui etiam sint inter factores numeri $A$, et quemvis horum factorum toties in $B, C, D$ etc. coniunctim occurrere debere, quoties in $A$. Hinc colligitur criterium, utrum numerus $B$ alium $A$ metiatur, necne. Illud eveniet, si $B$ neque alios factores primos, neque ullum pluries involvit, quam $A$; quarum conditionum si aliqua deficit, $B$ ipsum $A$ non metietur.

Facile hinc calculi combinationum auxilio derivari potest, si $A=a^\alpha b^\beta c^\gamma$ etc. designantibus ut supra $a, b, c$ etc. numeros primos diversos: $A$ habere
\[
(\alpha+1)(\beta+1)(\gamma+1)\ \mathrm{etc.}
\]
divisores diversos, inclusis etiam $1$ et $A$.

\Article{18}
Si igitur $A=a^\alpha b^\beta c^\gamma$ etc., $K=k^\kappa l^\lambda m^{\mu}$ etc., atque primi $a, b, c$ etc., $k, l, m$ etc. omnes diversi, patet $A$ et $K$ divisorem communem praeter $1$ non habere, sive inter se esse primos.

Pluribus numeris $A, B, C$ etc. propositis \emph{maxima omnibus communis mensura} ita determinatur. Resolvantur omnes in suos factores primos, atque ex his excerpantur ii, qui omnibus numeris $A, B, C$ etc. sunt communes (si tales non adsunt, nullus divisor erit omnibus communis). Tum quoties quisque horum factorum primorum in singulis $A, B, C$ etc. contineatur, sive \emph{quot dimensiones} in singulis $A, B, C$ etc. quisque habeat, adnotetur. Tandem singulis factoribus primis tribuantur dimensiones omnium quas in $A, B, C$ etc. habent minimae, componaturque productum ex iis, quod erit mensura communis quaesita.

Quando vero numerorum $A, B, C$ etc. \emph{minimus communis dividuus} desideratur, ita procedendum. Colligantur omnes numeri primi, qui numerorum $A, B, C$ etc. aliquem metiuntur, tribuatur cuivis dimensio omnium quas in numeris $A, B, C$ etc. habet maxima, sicque ex omnibus productum confletur, quod erit dividuus quaesitus.

\emph{Ex.} Sit $A=504=2^3 3^2 7$; $B=2880=2^6 3^2 5$; $C=864=2^5 3^3$. Pro inveniendo divisore communi maximo habentur factores primi $2, 3$, quibus dimensiones $3, 2$ tribuendi; unde fiet $=2^3 3^2=72$; dividuus vero communis minimus erit $2^6 3^3 5.7=60480$.

% Physical source leaf 022; retained native image authority.
\SourceLeaf{22}
Demonstrationes propter facilitatem omittimus. Ceterum quomodo haec problemata solvenda sint, quando numerorum $A, B, C$ etc. in factores resolutio non detur, ex elementis notum.

\Article{19}
\emph{Si numeri $a, b, c$ etc. ad alium $k$ sunt primi, etiam productum ex illis $abc$ etc. ad $k$ primum est.}

Quia enim nulli numerorum $a, b, c$ etc. factor primus cum $k$ est communis productumque $abc$ etc. alios factores primos habere nequit, quam qui sunt factores alicuius numerorum $a, b, c$ etc., productum $abc$ etc. etiam cum $k$ factorem primum communem non habebit. Quare ex art. praec. $k$ ad $abc$ etc. primus.

\emph{Si numeri $a, b, c$ etc. inter se sunt primi, aliumque $k$ singuli metiuntur: etiam productum ex illis numerum $k$ metietur.}

Hoc aeque facile ex artt. 17, 18 derivatur. Sit enim quicunque producti $abc$ etc. divisor primus $p$, quem contineat $\pi$ vicibus, manifestumque est, aliquem numerorum $a, b, c$ etc. eundem hunc divisorem $\pi$ vicibus continere debere. Quare etiam $k$, quem hic numerus metitur.\ $\pi$ vicibus divisorem $p$ continet. Similiter de reliquis producti $abc$ etc. divisoribus.

Hinc \emph{si duo numeri $m, n$ secundum plures modulos inter se primos $a, b, c$ etc. sunt congrui, etiam secundum productum ex his congrui erunt.} Quum enim $m-n$ per singulos $a, b, c$ etc. sit divisibilis, etiam per eorum productum dividi poterit.

Denique \emph{si $a$ ad $b$ primus et $ak$ per $b$ divisibilis, erit etiam $ak$ per $b$ divisibilis.} Namque quoniam $ak$ tam per $a$ quam per $b$ divisibilis, etiam per $ab$ dividi poterit, i. e. $\frac{ak}{ab}=\frac{k}{b}$ erit integer.

\Article{20}
\emph{Quando $A=a^\alpha b^\beta c^\gamma$ etc., designantibus $a, b, c$ etc. numeros primos inaequales, est potestas aliqua, puta $=k^n$: omnes exponentes $\alpha, \beta, \gamma$ etc. per $n$ erunt divisibiles.}

Numerus enim $k$ alios factores primos quam $a, b, c$ etc. non involvit. Contineat factorem $a$, $\alpha'$ vicibus, continebitque $k^n$ sive $A$ hunc factorem $n\alpha'$ vicibus; quare $n\alpha'=\alpha$, et $\frac{\alpha}{n}$ integer. Similiter $\frac{\beta}{n}$ etc. integros esse demonstratur.
